Consider an undirected graph that consists of $n$ nodes and $m$ edges. There are two types of events that can happen:

A new edge is created between nodes $a$ and $b$.
An existing edge between nodes $a$ and $b$ is removed.

Your task is to report the number of components after every event.

\vspace{0.3cm}\noindent\textbf{Input}

The first input line has three integers $n$, $m$ and $k$: the number of nodes, edges and events.
After this there are $m$ lines describing the edges. Each line has two integers $a$ and $b$: there is an edge between nodes $a$ and $b$. There is at most one edge between any pair of nodes.
Then there are $k$ lines describing the events. Each line has the form "$t$ $a$ $b$" where $t$ is 1 (create a new edge) or 2 (remove an edge). A new edge is always created between two nodes that do not already have an edge between them, and only existing edges can get removed.

\vspace{0.3cm}\noindent\textbf{Output}

Print $k+1$ integers: first the number of components before the first event, and after this the new number of components after each event.

\vspace{0.3cm}\noindent\textbf{Constraints}

\textbullet\ $2 \le n \le 10^5$
\textbullet\ $1 \le m,k \le 10^5$
\textbullet\ $1 \le a,b \le n$